\documentclass[10pt,a4paper]{amsart}
\usepackage{amsmath,amssymb,amsthm,graphicx}
% Retained minimal preamble for the Stacks Project chapter varieties.tex.
% Verbatim excerpt of the chapter preamble: Xy-pic setup (native preamble lines 14-18), the native theorem
% declarations with the mathematical macros (native preamble lines 154-217), and the project own no-op
% definitions of the reference/slogan/history display environments (native wrapper.tex lines 1-3) that the
% counted statement uses for its slogan. Copied mechanically; nothing added, removed or redefined.
% For commutative diagrams we use Xy-pic
\usepackage[all]{xy}

% We use 2cell for 2-commutative diagrams.
\xyoption{2cell}
% Theorem environments.
%
\theoremstyle{plain}
\newtheorem{theorem}[subsection]{Theorem}
\newtheorem{proposition}[subsection]{Proposition}
\newtheorem{lemma}[subsection]{Lemma}

\theoremstyle{definition}
\newtheorem{definition}[subsection]{Definition}
\newtheorem{example}[subsection]{Example}
\newtheorem{exercise}[subsection]{Exercise}
\newtheorem{situation}[subsection]{Situation}

\theoremstyle{remark}
\newtheorem{remark}[subsection]{Remark}
\newtheorem{remarks}[subsection]{Remarks}

\numberwithin{equation}{subsection}

% Macros
%
\def\lim{\mathop{\mathrm{lim}}\nolimits}
\def\colim{\mathop{\mathrm{colim}}\nolimits}
\def\Spec{\mathop{\mathrm{Spec}}}
\def\Hom{\mathop{\mathrm{Hom}}\nolimits}
\def\Ext{\mathop{\mathrm{Ext}}\nolimits}
\def\SheafHom{\mathop{\mathcal{H}\!\mathit{om}}\nolimits}
\def\SheafExt{\mathop{\mathcal{E}\!\mathit{xt}}\nolimits}
\def\Sch{\mathit{Sch}}
\def\Mor{\mathop{\mathrm{Mor}}\nolimits}
\def\Ob{\mathop{\mathrm{Ob}}\nolimits}
\def\Sh{\mathop{\mathit{Sh}}\nolimits}
\def\NL{\mathop{N\!L}\nolimits}
\def\CH{\mathop{\mathrm{CH}}\nolimits}
\def\proetale{{pro\text{-}\acute{e}tale}}
\def\etale{{\acute{e}tale}}
\def\QCoh{\mathit{QCoh}}
\def\Ker{\mathop{\mathrm{Ker}}}
\def\Im{\mathop{\mathrm{Im}}}
\def\Coker{\mathop{\mathrm{Coker}}}
\def\Coim{\mathop{\mathrm{Coim}}}

% Boxtimes
%
\DeclareMathSymbol{\boxtimes}{\mathbin}{AMSa}{"02}

%
% Macros for moduli stacks/spaces
%
\def\QCohstack{\mathcal{QC}\!\mathit{oh}}
\def\Cohstack{\mathcal{C}\!\mathit{oh}}
\def\Spacesstack{\mathcal{S}\!\mathit{paces}}
\def\Quotfunctor{\mathrm{Quot}}
\def\Hilbfunctor{\mathrm{Hilb}}
\def\Curvesstack{\mathcal{C}\!\mathit{urves}}
\def\Polarizedstack{\mathcal{P}\!\mathit{olarized}}
\def\Complexesstack{\mathcal{C}\!\mathit{omplexes}}
% \Pic is the operator that assigns to X its picard group, usage \Pic(X)
% \Picardstack_{X/B} denotes the Picard stack of X over B
% \Picardfunctor_{X/B} denotes the Picard functor of X over B
\def\Pic{\mathop{\mathrm{Pic}}\nolimits}
\def\Picardstack{\mathcal{P}\!\mathit{ic}}
\def\Picardfunctor{\mathrm{Pic}}
\def\Deformationcategory{\mathcal{D}\!\mathit{ef}}
\newenvironment{reference}{}{}
\newenvironment{slogan}{}{}
\newenvironment{history}{}{}
\usepackage{cleveref}
\crefname{originalreference}{original reference}{original references}
\Crefname{originalreference}{Original reference}{Original references}
\makeatletter
\newlabel{limits-proposition-affine}{{05YU}{source}{}{sourcepointer-1}{}}
\newlabel{limits-proposition-affine@cref}{{[originalreference][0][]05YU}{[1][0][]source}{}{}{}}
\makeatother
\crefname{originalreference}{original reference}{original references}
\Crefname{originalreference}{Original reference}{Original references}
\makeatletter
\newlabel{algebra-lemma-silly}{{00DS}{source}{}{sourcepointer-2}{}}
\newlabel{algebra-lemma-silly@cref}{{[originalreference][0][]00DS}{[1][0][]source}{}{}{}}
\makeatother
\crefname{originallemma}{original lemma}{original lemmas}
\Crefname{originallemma}{Original lemma}{Original lemmas}
\makeatletter
\newlabel{lemma-extra-silly}{{09NK}{source}{}{sourcepointer-3}{}}
\newlabel{lemma-extra-silly@cref}{{[originallemma][0][]09NK}{[1][0][]source}{}{}{}}
\makeatother
\crefname{originalreference}{original reference}{original references}
\Crefname{originalreference}{Original reference}{Original references}
\makeatletter
\newlabel{limits-lemma-affines-glued-in-closed-affine}{{09NL}{source}{}{sourcepointer-4}{}}
\newlabel{limits-lemma-affines-glued-in-closed-affine@cref}{{[originalreference][0][]09NL}{[1][0][]source}{}{}{}}
\makeatother
\title{A finite set of points whose traces on each irreducible component lie in an affine open lies in an affine open}
\author{The Stacks Project Authors}
\date{stacks-project:varieties.tex:local-snapshot-20261011; GNU FDL 1.2 or later}
\begin{document}
\maketitle
\noindent\textit{Editorial note.} All delivered mathematical text is the author's own, copied verbatim from the pinned Stacks Project chapter file varieties.tex; no statement or proof step was rewritten, reordered or removed.\par\medskip
\noindent\textit{Editorial note.} The counted span cites four cross-chapter labels, each recorded in assembly.external\_references with its canonical tag as a navigation pointer.\par\medskip
\noindent\textit{Editorial note.} The proof cites no bibliography item inside the counted span.\par\medskip
\noindent\textit{Allowed standard prerequisites (editorial).}\begin{itemize}
\item Affineness criterion via the limit presentation and prime avoidance. Stacks Project, Limits of Schemes, Proposition limits-proposition-affine (tag 01Z2); Algebra, Lemma algebra-lemma-silly (tag 00EL). The first is the criterion applied throughout and the second produces the standard opens avoiding the closed sets Delta and Delta'.
\end{itemize}\medskip
\ifdefined\hypertarget\hypertarget{sourcepointer-1}{}\fi
\noindent\textit{Source reference.} Original reference 05YU: sources/stacks-chow/tags.txt, source lines 7668--7668. Cross-chapter standard prerequisite cited by the author inside the counted span; it is used as a pointer and its own proof is not part of this item.\par\medskip
\ifdefined\hypertarget\hypertarget{sourcepointer-2}{}\fi
\noindent\textit{Source reference.} Original reference 00DS: sources/stacks-chow/tags.txt, source lines 508--508. Cross-chapter standard prerequisite cited by the author inside the counted span; it is used as a pointer and its own proof is not part of this item.\par\medskip
\ifdefined\hypertarget\hypertarget{sourcepointer-3}{}\fi
\noindent\textit{Source reference.} Original lemma 09NK: sources/stacks-chow/tags.txt, source lines 12218--12218. Cross-chapter standard prerequisite cited by the author inside the counted span; it is used as a pointer and its own proof is not part of this item.\par\medskip
\ifdefined\hypertarget\hypertarget{sourcepointer-4}{}\fi
\noindent\textit{Source reference.} Original reference 09NL: sources/stacks-chow/tags.txt, source lines 12219--12219. Cross-chapter standard prerequisite cited by the author inside the counted span; it is used as a pointer and its own proof is not part of this item.\par\medskip
\begin{lemma}
\label{lemma-finite-set-codim-1-points-in-affine-per-component}
Let $X$ be a scheme. Let $T \subset X$ be finite set of points. Assume
\begin{enumerate}
\item $X$ has finitely many irreducible components $Z_1, \ldots, Z_t$, and
\item $Z_i \cap T$ is contained in an affine open of the reduced
induced subscheme corresponding to $Z_i$.
\end{enumerate}
Then there exists an affine open subscheme of $X$ containing $T$.
\end{lemma}

\begin{proof}
Using Limits, Proposition \ref{limits-proposition-affine}
there is an immediate reduction to the case where $X$
is reduced. Details omitted. In the rest of the proof we
endow every closed subset of $X$ with the induced reduced closed
subscheme structure.

\medskip\noindent
We argue by induction that we can find an affine open
$U \subset Z_1 \cup \ldots \cup Z_r$ containing
$T \cap (Z_1 \cup \ldots \cup Z_r)$. For $r = 1$ this holds by assumption.
Say $r > 1$ and let $U \subset Z_1 \cup \ldots \cup Z_{r -  1}$
be an affine open containing $T \cap (Z_1 \cup \ldots \cup Z_{r - 1})$.
Let $V \subset X_r$ be an affine open containing $T \cap Z_r$ (exists by
assumption). Then $U \cap V$ contains
$T \cap ( Z_1 \cup \ldots \cup Z_{r - 1} ) \cap Z_r$.
Hence
$$
\Delta = (U \cap Z_r) \setminus (U \cap V)
$$
does not contain any element of $T$. Note that $\Delta$ is a closed
subset of $U$. By prime avoidance (Algebra, Lemma \ref{algebra-lemma-silly}),
we can find a standard open $U'$ of $U$ containing $T \cap U$ and avoiding
$\Delta$, i.e., $U' \cap Z_r \subset U \cap V$.
After replacing $U$ by $U'$ we may assume that $U \cap V$ is
closed in $U$.

\medskip\noindent
Using that by the same arguments as above also the set
$\Delta' = (U \cap (Z_1 \cup \ldots \cup Z_{r - 1})) \setminus (U \cap V)$
does not contain any element of $T$ we find a $h \in \mathcal{O}(V)$
such that $D(h) \subset V$ contains $T \cap V$ and such that
$U \cap D(h) \subset U \cap V$. Using that $U \cap V$ is closed in $U$
we can use Lemma \ref{lemma-extra-silly}
to find an element $g \in \mathcal{O}(U)$ whose restriction to $U \cap V$
equals the restriction of $h$ to $U \cap V$ and such that
$T \cap U \subset D(g)$. Then we can replace $U$ by $D(g)$
and $V$ by $D(h)$ to reach the situation where $U \cap V$ is
closed in both $U$ and $V$. In this case the scheme $U \cup V$ is
affine by Limits, Lemma \ref{limits-lemma-affines-glued-in-closed-affine}.
This proves the induction step and thereby the lemma.
\end{proof}

\end{document}
